\subsection*{Spettro Continuo}
Negli spazi di Hilbert a infinite dimensioni, è possibile che la base sia continua invece che discreta. In questo caso, gli autovalori dell'operatore non sono numerabili e l'espansione di uno stato in termini di autostati diventa un'integrale:
\[
    |\psi\ket = \int \, \bra \xi | \psi \ket |\xi \ket d\xi
\]
dove $\bra \xi | \psi \ket$ è il proiettore con variazione continua.
Risulta che
\[
    \bra \xi | \xi' \ket = c_0\delta(\xi - \xi')
\]
Inoltre un autovettore in uno spazio di questo tipo ha questo aspetto:
\[
    \hat X |\xi\ket = \xi |\xi\ket
\]
E l'espansione dell'identità diventa
\[
    \hat Id = \int |\xi\ket \bra \xi | d\xi
\]

\subsubsection*{Un esempio}
Consideriamo gli operatori posizione, che funzionano in questo modo:

\begin{align*}
    |\vec x\ket = |x,y,z\ket \\
    \hat{X}_j |\vec x\ket = \vec x |\vec x\ket \\
    \hat{Y} |\vec x\ket = \vec y |\vec x\ket \\
    \hat{Z} |\vec x\ket = \vec z |\vec x\ket
\end{align*}

\subsubsection*{Altro esempio: operatore di traslazione}
Chiamiamo $\tau$ l'operatore che si comporta in questo modo:
\[
    \tau (\vec {dx}) |\vec x\ket = |\vec x + \vec {dx}\ket
\]
Al chè si può usare il solito trucco:
\[
    \tau (\vec {dx}) \hat{Id} |\psi\ket = \rnint |\vec x + \vec {dx}\ket \bra \vec x + \vec {dx} | \psi\ket d^nx
\]
Inoltre serve che $\tau$ rispetti altre proprietà:
\begin{align*}
    \bra\psi|\tau^\dag (\vec{dx}) \tau(\vec{dx})|\psi\ket = \hat{Id} \\
    \tau(\vec{dx}_1) \tau(\vec{dx}_2) = \tau(\vec{dx}_1 + \vec{dx}_2) \\
    \tau(\vec{0}) = \hat{Id}
\end{align*}

Una valida implementazione di $\tau$ è la seguente:
\[
    \tau (\vec {dx}) = \hat{Id} - i \hat k \cdot \vec{dx}
\]
Si può verificare che l'ultima espressione è sufficiente per rispettare le proprietà richieste, a patto che $\vec{k}$ sia un operatore hermitiano ($\hat{k} = \hat{k} ^ \dag$). \\
A questo punto si ha anche
\[
    \tau^\dag (\vec{dx}) = \hat{Id} + i \hat k \cdot \vec{dx}
\]
Si può provare a calcolare il commutatore tra posizione e traslazione:
\begin{align*}
    [\hat{X}_j, \tau(\vec{dx})] |\vec x\ket &= \hat{X}_j \tau(\vec{dx}) |\vec x\ket - \tau(\vec{dx}) \hat{X}_j |\vec x\ket \\
    &= \hat{X}_j |\vec x + \vec{dx}\ket - \tau(\vec{dx}) \vec x |\vec x\ket \\
    &= (\vec{dx}) |\vec x + \vec{dx}\ket
\end{align*}
per ogni operatore posizione al variare di $j$. \\
Risulta infine:
\[
    [\hat{X}_i, \hat{k}_j] = \delta_{ij} \hat{Id}
\]
ma ricordando che $[\hat{X}_j, \hat{p}_l] = i \hslash \delta_{jl} \hat{Id}$, si ottiene che
\[
    \hat{p}_l = i\hslash \hat{k}_l
\]
Però ora è possibile indurre questo ragionamento: per avere una traslazione finita, è possibile ripetere molte volte una traslazione infinitesima. Quindi:
\begin{align*}
    \tau(\vec{x}) &= \lim_{N \to \infty} \left[ \tau\left(\frac{\vec{x}}{N}\right) \right]^N \\
    &= \lim_{N \to \infty} \left[ \hat{Id} - i \hat{k} \cdot \frac{\vec{x}}{N} \right]^N \\
    &= e^{-i \hat{k} \cdot \vec{x}} \\
    &= e^{-\frac{i}{\hslash} \hat{p} \cdot \vec{x}}
\end{align*}
Inoltre i $\hat p$ devono commutare, così da avere traslazioni indipendenti nelle tre direzioni spaziali:
\[
    [\hat{p}_j, \hat{p}_l] = 0
\]
E, siccome $\hat p |\vec p\ket = \vec p |\vec p\ket$, si ha
\[
    \tau(\vec x)|\vec p\ket = e^{-\frac{i}{\hslash} \hat p \cdot \vec x} |\vec p\ket = e^{-\frac{i}{\hslash} \vec p \cdot \vec x} |\vec p\ket
\]

All'inizio del prossimo capitolo verrà spiegato il significato di funzione di operatore.

\subsubsection*{Rappresentazione nella base dei momenti}
Proviamo ad espandere $\tau(\vec{dx})|\psi\ket$:
\begin{gather*}
    \tau(\vec{dx})|\psi\ket = \tau(\vec{dx}) \rnint |\vec x \ket \bra \vec x | \psi\ket d^nx =\\
    = \rnint |\vec x + \vec{dx} \ket \bra \vec x || \psi\ket d^nx =  \rnint |\vec x \ket \bra \vec x - \vec{dx} | \psi\ket d^nx = \\
    = \rnint |\vec x \ket \psi(\vec x - \vec{dx}) d^nx = \rnint |\vec x \ket \left[ \psi(\vec x) - \vec{dx} \cdot \Grad \psi(\vec x) \right] d^nx = \\
    = \rnint |\vec x \ket \left[ \bra \vec x | \psi \ket - \vec{dx} \cdot \Grad \bra \vec x | \psi \ket \right] d^nx
    = \rnint \ketbrabra {\vec x} - |\vec x \ket \bra \vec x |  \vec{dx} \cdot \Grad |\psi\ket d^nx = \\
    = \left[\hat {Id} - \rnint \vec{dx} \cdot \hat \Grad d^nx \right] \Ket \psi
\end{gather*}
Ma ricordando la definizione di $\tau$ risulta che:
\begin{gather*}
    \tau(\vec{dx})\Ket \psi = \l \cancel {\hat {Id}} - i\frac{\hat p}{\h} \cdot \vec{dx}\r\Ket \psi = \left[\cancel {\hat {Id}} - \rnint \vec{dx} \cdot \hat \Grad d^nx \right] \Ket \psi \\
    \Bra {\vec x} \hat p \Ket \psi = \Bra {\vec x} -i\h\hat\Grad\Ket\psi = -i\h\hat\Grad \braket{\vec x}{\psi}
\end{gather*}

Quest'ultima equazione esprime l'azione dell'operatore momento ($\hat p$) quando lo stato è espresso nella base delle posizioni: in particolare ne è il gradiente!

\subsubsection*{Autofunzioni del momento nella base delle posizioni}
\begin{gather*}
    \hat p |\vec p\ket = \vec p |\vec p\ket \\
    \bra \vec x | \hat p | \vec p\ket = \vec p \bra \vec x | \vec p\ket = -i \hslash \Grad \bra \vec x | \vec p\ket
\end{gather*}
\newcommand{\braketxp}{\braket {\vec x} {\vec p}}
\newcommand{\braketpx}{\braket {\vec p} {\vec x}}
Dove $\Ket {\vec p}$ è un'autostato dell'operatore momento. Quest'ultima è un'equazione differenziale in $\braketxp$, le cui soluzioni improprie sono del tipo:
\[
    \braketxp = A e^{\frac{i}{\hslash} \vec p \cdot \vec x}
\]
Per trovare la $A$ si può procedere come segue:
\begin{gather*}
    \bra \vec x' | \vec x'' \ket = \delta(\vec x' - \vec x'') 
    = \rnint \bra \vec x' | \vec p \ket \bra \vec p | \vec x'' \ket d^np = \\
    = |A|^2 \rnint e^{\frac{i}{\hslash} \vec p \cdot (\vec x' - \vec x'')} d^np
    = |A|^2 (2 \pi \hslash)^{n/2} \delta(\vec x' - \vec x'')
\end{gather*}
Al chè $A = \frac{1}{(2 \pi \hslash)^{n/2}}$.\\
Assemblando quanto trovato fin'ora risulta:
\begin{gather*}
    \bra \vec x | \psi \ket = \psi(\vec x) = \\
    = \rnint \braketxp \braket {\vec p}{\psi} d^np = \\
    = \frac{1}{(2 \pi \hslash)^{n/2}} \rnint e^{\frac{i}{\hslash} \vec p \cdot \vec x} \psi(\vec p) d^np
\end{gather*}
Che coincide con la definizione della trasformata di Fourier! E anche: siccome invertire i termini del prodotto scalare coincide col fare il complesso coniugato risulta anche:

Assemblando quanto trovato fin'ora risulta:
\begin{gather*}
    \bra \vec p | \psi \ket = \psi(\vec p) = \\
    = \rnint \braketpx \braket {\vec x}{\psi} d^nx = \\
    = \frac{1}{(2 \pi \hslash)^{n/2}} \rnint e^{-\frac{i}{\hslash} \vec p \cdot \vec x} \psi(\vec x) d^nx
\end{gather*}
Ovvero l'antitrasformata!
